% Identifiability of the parameters selected by the other criteria
% (tables: reproducibility/figures_tables/identif_criteria_tables.py; profiles: scripts/run_identifiability_other_criteria.sh).

\snote{identcrit}

\emph{Question.} Brun's collinearity index, profile likelihood and our conditioning tests all aim, in different ways, at subsets of identifiable parameters. Brun's index and our conditioning tests judge identifiability from the sensitivity matrix at the starting point, that is, locally and to first order; profile likelihood judges it from the profiles of the full model; Morris screening and pivoted QR do not include an explicit redundancy test against compensation. The main text shows that the subsets of these criteria differ in whether they reproduce the full model. This note asks whether they also differ in the practical identifiability of the parameters they select, measured in the same way for every criterion, and whether any criterion other than ours obtains subsets that are at the same time identifiable and sufficient.

\emph{Design.} For every problem and criterion, the subset of the criterion (selected on the training outputs, with the same number $k=|S|$ of biological parameters as our method; Methods, \emph{Comparison with other selection criteria}) was profiled with exactly the procedure used for our subsets (\sref{identif}): the reduced problem has the selected and the calibration parameters free and all other parameters at $\theta_0$; it is fitted to all measurements from $\theta_0$; for every selected parameter, $\log|\theta_j|$ is fixed at ten displacements per side, from 0.05 to 4.6, and the other free parameters are re-estimated; points outside the bounds $[|\theta_0|/10,10|\theta_0|]$ are not used; if a point improves the optimum the problem is re-fitted from it and the profile repeated (at most three rounds); and a parameter is identifiable if $\Delta\chi^2$ exceeds 3.84 on both sides. Where a criterion chose exactly our subset, our profiles were used (10 problems for Brun, 6 for Morris, 2 for profile likelihood and 1 for pivoted QR). The profiles were computed on a 128-core server (\texttt{scripts/run\_identifiability\_other\_criteria.sh}, which calls \texttt{analysis/prepare\_criteria\_subsets.py} and \texttt{analysis/identifiability\_profiles.py}); the tables are produced by \texttt{figures\_tables/identif\_criteria\_tables.py}.

The comparison is made on the 27 problems in which every criterion has a subset: Blasi has no training/test split on which the other criteria could select, Chen was not included in the selections of the other criteria (Methods), Raimundez has no profile-likelihood subset, and Bertozzi is excluded as in the comparison of admissibility. For each criterion we count (i) the selected parameters that are identifiable, (ii) the problems in which every selected parameter is identifiable, and (iii) the problems in which every selected parameter is identifiable \emph{and} the subset is admissible at $\pm$5\% (Table~2 of the main text), and we compare (ii) and (iii) with our subsets by the exact McNemar test over problems, with Holm adjustment over the four existing criteria. Our subsets selected on the training outputs are reported as well, on the 26 of these problems in which that selection gave an admissible subset.

\emph{Results.} Supplementary Table~\ref{tab:identifcrit} gives the summary and Supplementary Table~\ref{tab:identifcritprob} the result for every problem and criterion.

\begin{itemize}\itemsep0pt
\item \emph{Identifiability alone separates only the criteria without a redundancy test.} 86\% of our selected parameters were identifiable (57 of 66), 85\% of Brun's (56), 86\% of those of profile likelihood (57), 80\% of Morris's (53) and 76\% of those of pivoted QR (50). The number of problems in which every selected parameter is identifiable was 21 for our subsets, 20 for Brun, 22 for profile likelihood, 18 for Morris and 14 for pivoted QR; none of these differences from our method is significant after adjustment (Holm-adjusted $P=1$, 1, 0.75 and 0.16).
\item \emph{The same identifiability is reached with different parameters.} Brun's subsets share two thirds of our selected parameters (46 of 70 over its 29 problems) and those of profile likelihood about two fifths (26 of 68), yet both are as identifiable as ours: a model usually contains several groups of parameters that the data determine, and the criteria find different ones.
\item \emph{Requiring identifiability and sufficiency together separates the criteria.} Every selected parameter was identifiable and the subset reproduced the full model in 21 of the 27 problems with our subsets, 18 with Brun's, 14 with Morris's, 14 with those of profile likelihood and 7 with those of pivoted QR (Holm-adjusted McNemar $P=0.38$, 0.047, 0.078 and 0.002; problems in which only our subsets have both properties against problems in which only the other criterion has them: 4 against 1, 7 against 0, 8 against 1 and 15 against 1). Selected on the training outputs, our subsets had both properties in 16 of 26 problems.
\end{itemize}

\emph{Interpretation.} Profile likelihood chooses, by construction, the parameters with the narrowest confidence intervals, and its subsets are therefore as identifiable as ours; but in a third of the problems they do not reproduce the full model when all parameters change (18 of 27 admissible). Brun's collinearity index, which judges identifiability only locally and to first order, gives subsets that are in practice as identifiable as ours and almost as often sufficient, consistent with the close similarity between its ranking and ours (Discussion). Morris screening and pivoted QR give less identifiable subsets that are also less often sufficient. The property that distinguishes the subsets of our method is thus not identifiability by itself, which several criteria achieve, but the combination of identifiability with the verified ability to reproduce the full model, which is what the method tests by re-estimation.

\emph{Computing time.} The profiles of the subsets of the four criteria took 41~h of processor time (Brun 4.7~h, Morris 13.1~h, profile likelihood 8.6~h, pivoted QR 14.4~h), most of it for Isensee, whose profile points take up to an hour each; on the 128-core server the whole analysis took about 2~h of elapsed time.

\emph{Limitations.} The comparison uses the profiles of the reduced problem with the non-selected parameters at $\theta_0$, as for our subsets; it does not assess identifiability in the full model. The subsets of the other criteria were selected on the training outputs, whereas the profiles use all measurements, as for our subsets. Pivoted QR, Morris screening and profile likelihood were applied as they are usually implemented (\sref{morrisbrun}); other implementations could select different parameters.

{\scriptsize\setlength{\tabcolsep}{3pt}
\begin{longtable}{@{}lrrrrrrr@{}}
\caption{Practical identifiability of the parameters selected by every criterion, with the profile-likelihood procedure of our subsets (reduced problem, all data, $\Delta\chi^2=3.84$ on both sides), over the 27 problems in which every criterion has a subset. Identifiable: selected parameters whose 95\% profile-likelihood interval is bounded on both sides within the bounds of the study. All identifiable: problems in which every selected parameter is identifiable. Admissible: the subset reproduces the full model at $\pm$5\% (Table~2 of the main text). Both: all parameters identifiable \emph{and} the subset admissible. $P$: exact McNemar test against our subsets (all outputs) over problems, for ``all identifiable'' and for ``both'', Holm-adjusted over the four existing criteria (for our subsets selected on the training outputs, unadjusted). Bertozzi is excluded, as in the comparison of admissibility. Same as ours: problems in which the criterion chose exactly our subset, whose profiles are then ours. Our subsets selected on the training outputs are shown only for the problems in which that selection gave an admissible subset.}\label{tab:identifcrit}\\
\toprule
Criterion & Problems & Identifiable & All identifiable & Admissible & Both & $P$ (all id.; both) & Same as ours \\
\midrule
\endfirsthead
\toprule
Criterion & Problems & Identifiable & All identifiable & Admissible & Both & $P$ (all id.; both) & Same as ours \\
\midrule
\endhead
Ours (all outputs) & 27 & 57/66 (86\%) & 21 & 27 & 21 & -- & -- \\
Ours (training) & 26 & 50/61 (82\%) & 19 & 23 & 16 & 0.5; 0.062 & 9 \\
Brun & 27 & 56/66 (85\%) & 20 & 25 & 18 & 1; 0.38 & 9 \\
Morris & 27 & 53/66 (80\%) & 18 & 22 & 14 & 0.75; 0.047 & 6 \\
Profile likelihood & 27 & 57/66 (86\%) & 22 & 18 & 14 & 1; 0.078 & 2 \\
Pivoted QR & 27 & 50/66 (76\%) & 14 & 16 & 7 & 0.16; 0.0021 & 1 \\
\bottomrule
\end{longtable}

}

{\small
\begin{longtable}{@{}lcccccc@{}}
\caption{Identifiable parameters / selected parameters for every problem and criterion (profile likelihood of the reduced problem). $=$, same subset as ours (profiles shared); $^{\ast}$, the subset is admissible at $\pm$5\%; --, no subset. Problems below the line are not common to all criteria and are not included in the summary of Supplementary Table~\ref{tab:identifcrit}.}\label{tab:identifcritprob}\\
\toprule
Problem & Ours (all outputs) & Ours (training) & Brun & Morris & Profile likelihood & Pivoted QR \\
\midrule
\endfirsthead
\toprule
Problem & Ours (all outputs) & Ours (training) & Brun & Morris & Profile likelihood & Pivoted QR \\
\midrule
\endhead
Alkan & 2/2$^{\ast}$ & 2/2$^{\ast}$ & 2/2$^{\ast}$ & 2/2$^{\ast}$ & 2/2$^{\ast}$ & 2/2 \\
Armistead & 2/2$^{\ast}$ & $=$ 2/2$^{\ast}$ & $=$ 2/2$^{\ast}$ & $=$ 2/2$^{\ast}$ & 2/2 & 2/2 \\
Bachmann & 2/2$^{\ast}$ & 2/2$^{\ast}$ & 2/2$^{\ast}$ & $=$ 2/2$^{\ast}$ & 2/2$^{\ast}$ & 1/2$^{\ast}$ \\
Boehm & 2/2$^{\ast}$ & $=$ 2/2$^{\ast}$ & $=$ 2/2$^{\ast}$ & $=$ 2/2$^{\ast}$ & 2/2 & 1/2 \\
Borghans & 2/2$^{\ast}$ & 2/2$^{\ast}$ & 2/2$^{\ast}$ & 2/2$^{\ast}$ & 2/2$^{\ast}$ & 2/2$^{\ast}$ \\
Brannmark & 2/2$^{\ast}$ & 2/2$^{\ast}$ & 2/2$^{\ast}$ & 0/2$^{\ast}$ & 2/2$^{\ast}$ & 1/2$^{\ast}$ \\
Bruno & 5/5$^{\ast}$ & 5/5 & 5/5 & 5/5 & 5/5 & 5/5 \\
Elowitz & 2/2$^{\ast}$ & $=$ 2/2$^{\ast}$ & $=$ 2/2$^{\ast}$ & 2/2 & 2/2 & 2/2 \\
Fiedler & 2/2$^{\ast}$ & 1/2$^{\ast}$ & 1/2$^{\ast}$ & 1/2$^{\ast}$ & 2/2$^{\ast}$ & 1/2$^{\ast}$ \\
Fujita & 0/2$^{\ast}$ & $=$ 0/2$^{\ast}$ & $=$ 0/2$^{\ast}$ & 1/2$^{\ast}$ & 2/2$^{\ast}$ & 2/2$^{\ast}$ \\
Giordano & 0/2$^{\ast}$ & 0/2$^{\ast}$ & 0/2$^{\ast}$ & 0/2$^{\ast}$ & 0/2 & 0/2 \\
Isensee & 2/2$^{\ast}$ & $=$ 2/2$^{\ast}$ & 1/2$^{\ast}$ & 2/2$^{\ast}$ & 2/2$^{\ast}$ & 2/2$^{\ast}$ \\
Laske & 2/2$^{\ast}$ & 2/2$^{\ast}$ & 2/2$^{\ast}$ & 2/2$^{\ast}$ & 2/2$^{\ast}$ & 2/2$^{\ast}$ \\
Liu & 2/2$^{\ast}$ & 2/2$^{\ast}$ & 2/2$^{\ast}$ & 2/2$^{\ast}$ & 2/2$^{\ast}$ & 2/2$^{\ast}$ \\
Lucarelli & 10/10$^{\ast}$ & 8/8 & 10/10$^{\ast}$ & 10/10 & 10/10 & 10/10 \\
Okuonghae & 1/2$^{\ast}$ & 0/2$^{\ast}$ & 1/2$^{\ast}$ & 0/2$^{\ast}$ & 0/2$^{\ast}$ & 1/2$^{\ast}$ \\
Oliveira & 1/2$^{\ast}$ & 0/2$^{\ast}$ & 1/2$^{\ast}$ & 1/2$^{\ast}$ & 1/2$^{\ast}$ & 0/2$^{\ast}$ \\
Perelson & 0/2$^{\ast}$ & -- & $=$ 0/2$^{\ast}$ & $=$ 0/2$^{\ast}$ & $=$ 0/2$^{\ast}$ & $=$ 0/2$^{\ast}$ \\
Rahman & 2/2$^{\ast}$ & 2/2$^{\ast}$ & 2/2$^{\ast}$ & 2/2$^{\ast}$ & 2/2$^{\ast}$ & 1/2$^{\ast}$ \\
Raia & 2/2$^{\ast}$ & 1/2$^{\ast}$ & $=$ 2/2$^{\ast}$ & 2/2$^{\ast}$ & 2/2$^{\ast}$ & 2/2$^{\ast}$ \\
Salazar-Cavazos & 2/2$^{\ast}$ & $=$ 2/2$^{\ast}$ & $=$ 2/2$^{\ast}$ & $=$ 2/2$^{\ast}$ & 2/2$^{\ast}$ & 2/2 \\
Schwen & 2/2$^{\ast}$ & $=$ 2/2$^{\ast}$ & 2/2$^{\ast}$ & $=$ 2/2$^{\ast}$ & $=$ 2/2$^{\ast}$ & 2/2$^{\ast}$ \\
Smith & 1/2$^{\ast}$ & $=$ 1/2$^{\ast}$ & 2/2$^{\ast}$ & 1/2$^{\ast}$ & 0/2$^{\ast}$ & 1/2$^{\ast}$ \\
Sneyd & 2/2$^{\ast}$ & $=$ 2/2$^{\ast}$ & $=$ 2/2$^{\ast}$ & 2/2 & 2/2 & 1/2 \\
Weber & 2/2$^{\ast}$ & 2/2$^{\ast}$ & $=$ 2/2$^{\ast}$ & 2/2$^{\ast}$ & 2/2$^{\ast}$ & 1/2$^{\ast}$ \\
Zhao & 3/3$^{\ast}$ & 2/2 & 3/3 & 2/3 & 3/3 & 3/3 \\
Zheng & 2/2$^{\ast}$ & 2/2$^{\ast}$ & 2/2$^{\ast}$ & 2/2$^{\ast}$ & 2/2 & 1/2 \\
\midrule
Bertozzi & 2/2 & -- & 2/2 & 1/2 & 2/2 & 2/2 \\
Blasi & 2/2 & $=$ 2/2 & -- & -- & -- & -- \\
Chen & 2/2 & $=$ 2/2 & -- & -- & -- & -- \\
Raimundez & 2/2$^{\ast}$ & $=$ 2/2$^{\ast}$ & $=$ 2/2$^{\ast}$ & 2/2$^{\ast}$ & -- & 2/2$^{\ast}$ \\
\bottomrule
\end{longtable}

}
